\chapter{«Veinte segundos»}
\label{sec:1}
Un oficial de inteligencia, al describir su papel en una cadena de selección de objetivos construida en torno a listas de muerte generadas por máquina, le contó a un periodista que dedicaba veinte segundos a cada nombre, lo justo para comprobar que el objetivo era un hombre: «No se tarda mucho en saber si alguien tiene voz de hombre o de mujer». Añadió que, como ser humano, su valor añadido era cero, salvo como sello de aprobación. «Ahorraba mucho tiempo» \autocite{abraham2024lavender}. Según su relato, a eso se reducía el «humano en el bucle» bajo el ritmo operativo.\footnote{La investigación sobre el sesgo de automatización predice esta delegación: cuando otras exigencias compiten por su atención, las personas a las que se da la recomendación de una máquina la comprueban menos y la acatan más \autocite{parasuraman2010complacency}.}

Una interpretación es que el oficial es un sujeto al que se le dio demasiado poco tiempo, y que la solución es gobernarlo mejor —más tiempo, un segundo revisor, más conciencia instalada en el bucle—. A través de esta lente de la \emph{sujeción social}, la persona es un sujeto que está fuera de una máquina y la maneja como una herramienta: un conductor y un coche, un obrero y un torno. La persona tiene un papel, un nombre, responsabilidades y una conciencia a la que uno puede dirigirse. Una lente alternativa es la \emph{servidumbre maquínica}, en la que la persona es un componente de la máquina, ensamblado con sus otras piezas en un sistema que las regula juntas; la máquina usa a la persona como uno de sus órganos, y lo que la persona aporta es una función realizada al ritmo que fija el sistema.\footnote{El nombre que Doctorow da a la posición del oficial es el centauro inverso: no una persona asistida por una máquina, sino un componente de la máquina hecho de una persona \autocite{doctorow2026reversecentaur}.} Las dos no se excluyen mutuamente; una persona puede tener un papel en un plano mientras es una pieza en el otro \autocite[Meseta 13, «The Apparatus of Capture»: hay servidumbre cuando los propios seres humanos son piezas constitutivas de una máquina que componen entre ellos y con otras cosas, bajo el control y la dirección de una unidad superior, y sujeción cuando la unidad superior constituye al ser humano como sujeto vinculado a un nuevo objeto exterior, que puede ser un animal, una herramienta o incluso una máquina]{deleuze1987plateaus}, y el oficial de los veinte segundos tenía las dos cosas. Como pieza, era el lugar donde la cadena producía su sello, y veinte segundos era el tiempo de ciclo de esa pieza. Como sujeto, se le interpelaba todo el tiempo: le describió a un periodista la disonancia entre las noticias de muertes civiles y la presión en el trabajo para conseguir más objetivos. E incluso cuando la disonancia llegó a su conciencia, siguió adelante. Se hizo notar y no decidió nada;\footnote{O, al menos, nada que impidiera a las Fuerzas de Defensa de Israel matar a cien mil personas en Gaza entre octubre de 2023 y octubre de 2025 \autocite{gomezugarte2025uncertainty,acosta2026gaza,economist2025gaza}. Y el relato de ese oficial en concreto lo publicó +972 Magazine en abril de 2024. Nadie puede saber si la decisión de ese oficial de hablar con un periodista evitó una muerte más, cincuenta mil muertes más o ninguna. ¿Y qué habría pasado si el oficial se hubiera negado tras dos semanas en el puesto —digamos, hipotéticamente, en noviembre de 2023—? Hablamos de un oficial de inteligencia que usaba Lavender para elaborar listas de muerte. Su entorno de trabajo diario probablemente se parecía más a un trabajo de oficina civil que a una película de acción. ¿Debería haber buscado reunirse una y otra vez con sus superiores, subiendo por la cadena de mando, mientras se negaba a aprobar objetivos para ataques aéreos de los que razonablemente cabía esperar muertes civiles? Podría decirse que eso habría sido lo ético. En ese contrafactual, ¿habría habido menos muertos?} el juicio de que la política era desproporcionada llegó, dijo, en retrospectiva.

Así que lo que faltaba no era una conciencia; la tenía. Una conciencia funciona haciendo que una decisión cueste algo, pero aquí el balance de costos salía al revés: «Nos presionaban constantemente: “Tráigannos más objetivos”. Nos gritaban de verdad» \autocite{abraham2024lavender}. Era un componente con escrúpulos, pero el ritmo se mantuvo: los escrúpulos no cambiaron nada de lo que aportaba, y se convirtió en un sello. Un componente con escrúpulos que no sostiene el ritmo se convierte en el «cuello de botella humano» sobre cuya eliminación escribió un libro un excomandante de la unidad de inteligencia de señales de Israel \autocite{ys2021humanmachine}. Lo que hacía del oficial un componente era cómo estaba dispuesto todo, no su actitud. Un componente puede seguir adelante a regañadientes o con entusiasmo: el funcionario irreflexivo que Arendt vio en el juicio de Eichmann \autocite{arendt1963eichmann} era, a juzgar por lo que había dicho en Argentina cuatro años antes, sin nada en juego, una actuación que montó para Jerusalén \autocite{stangneth2014eichmann}. Lo que la cadena dejaba aportar al oficial era una operación fija a un ritmo fijo: veía cada caso tal como la cadena lo describía, de uno en uno, en el tiempo que la cadena le daba. Si algo podría ocupar esa posición y no ser un componente es la pregunta de la parte IV.

Una investigación periodística describió tres sistemas que funcionaban juntos en Gaza: uno que generaba objetivos de edificios e infraestructuras, y según las fuentes la mayoría de sus sugerencias eran viviendas particulares; un segundo que puntuaba a individuos según su presunta afiliación militante y producía listas que, según se dijo, llegaron a marcar a hasta treinta y siete mil personas, y que internamente se sabía que se equivocaban más o menos una vez de cada diez\footnote{El ejército israelí cuestionó ese relato. Los relatos de Minab y de Maven que siguen se basan en investigaciones periodísticas que las partes también cuestionan. Los tomo tal como se publicaron.}; y un tercero que rastreaba cuándo las personas señaladas volvían a casa, para que los ataques aéreos pudieran programarse para cuando estuvieran con sus familias.

\runin{Maven: la versión estadounidense, ahora con aceleración por IA generativa}

El bombardeo de la escuela de Minab mató a 156 personas, 120 de ellas escolares \autocite[párr.~92]{UNFFMIran2026}.\footnote{Las cifras varían según la fuente y la fecha, desde 156 (el balance final del fiscal de Minab, que también da Amnistía Internacional) hasta 168 (el alcalde de la ciudad, en los primeros días). La misión de la ONU da las cifras del fiscal y además recoge información independiente de 157 muertos, 123 de ellos niños de 13 años o menos \autocite[párrs.~92–93]{UNFFMIran2026}.} La función civil de la escuela era visible en imágenes de satélite de fuentes abiertas, y la codificación del objetivo todavía llevaba una designación heredada de una instalación contigua de la Guardia Revolucionaria. Una misión de investigación de la ONU señaló que había motivos razonables para creer que el bombardeo constituyó el crimen de guerra de convertir un bien civil en objeto de un ataque \autocite{ohchr2026minabpress}, y que la escuela era claramente identificable sin ninguna información que indicara un uso militar \autocite[párrs.~81–86, 124–126]{UNFFMIran2026}. Concluyó que Estados Unidos incumplió su obligación de hacer todo lo factible para verificar el objetivo, contando la actualización de los datos de selección de objetivos entre las precauciones factibles, y situó el incumplimiento más allá de la negligencia \autocite[párrs.~124–125]{UNFFMIran2026}. Amnistía Internacional llegó a una conclusión paralela, en términos que siguen el deber de verificar del artículo 57(2)(a)(i) del Protocolo Adicional I: que las fuerzas pueden haber confiado en inteligencia desactualizada y no haber hecho todo lo factible para verificar el objetivo \autocite{Amnesty2026Minab}.

Ocurrió la mañana del 28 de febrero, el primer día de la campaña de Irán \autocite{Amnesty2026Minab,crs2026iranwar}. El día anterior, el secretario había anunciado que el proveedor (Anthropic) cuyo modelo (Claude) estaba dentro de la plataforma de selección de objetivos (Maven Smart System) sería designado riesgo para la cadena de suministro por negarse a licenciar Claude para todo uso lícito \autocite{npr2026openai}. Maven Smart System había devuelto la escuela como objetivo recomendado para el primer día \autocite{livingstone2026maven,gizmodo2026minab}. Ningún relato publicado dice qué papel tuvo Claude en esa recomendación, y una persona conocedora del trabajo del proveedor dijo a NBC News que el modelo no daba directamente asesoramiento sobre objetivos \autocite{nbc2026targeting}.

Una revisión interna del Pentágono, cuyas conclusiones completas siguen muy guardadas, habría concluido, según se informa, que la preparación de objetivos que en conflictos anteriores tardaba horas se condensó en minutos, que el personal dedicado a mitigar daños a civiles cayó aproximadamente un noventa por ciento, hasta menos de veinte personas, que el equipo del Mando Central pasó de diez a uno, y que ningún miembro de él revisó el emplazamiento de Minab —un resultado que los funcionarios atribuyeron tanto a los recortes como a la decisión de no implicar al equipo en la planificación y las operaciones— \autocite{Bloomberg2026KillChain}. La investigación formal se había abierto en marzo, después de que una evaluación preliminar concluyera que Estados Unidos era probablemente responsable \autocite{NPR2026MinabProbe}.

En la campaña de Irán de 2026, una plataforma de selección de objetivos construida por un contratista, con el modelo de uso general de un proveedor comercial dentro, usada para clasificar objetivos por importancia estratégica y redactar una justificación jurídica para cada ataque \autocite{livingstone2026maven}, ayudó a atacar más de mil objetivos en las primeras veinticuatro horas, unas diez veces lo que era posible antes de ella \autocite{csis2026maven}, y trece mil para el día treinta y ocho \autocite{breakingdefense2026maven}. El Pentágono no ha publicado ningún tiempo de revisión por objetivo para ella.\footnote{Maven es un nombre de programa que abarca más de un sistema: el programa del Pentágono que sigue llamándose Maven, y Maven Smart System, la plataforma de software que surgió de él \autocite{csis2026maven}. El programa empezó en 2017 como un proyecto para etiquetar objetos en el video que los drones recogían durante la campaña contra el Estado Islámico \autocite{weisgerber2017maven,fang2019maven}. El sistema de 2026 fusiona sistemas militares y fuentes de inteligencia en una sola interfaz que se usa para encontrar y atacar objetivos y para dirigir operaciones \autocite{defensescoop2026transition}. La plataforma es de Palantir: el director digital y de IA del Pentágono la llama «el producto de software como servicio de Palantir» \autocite{defensescoop2026transition}. La posición pública de Palantir, citada de segunda mano, es que decidir sobre la vida o la muerte no es su papel y que la responsabilidad siempre recae en la organización militar \autocite{lee2026palantir}, y sostiene que la inteligencia de base era responsabilidad del gobierno y que no hay pruebas de que Maven funcionara mal ni de que causara el bombardeo de Minab \autocite{Bloomberg2026KillChain}. Cómo pasó el programa de una oficina y un contratista a los siguientes se cuenta en el apéndice~\ref{sec:C.1}.}

Las personas que figuraban en las listas de ataques aéreos israelíes y estadounidenses no tenían manera de oponerse a figurar en esas listas: las cadenas no se construyeron para recabar ese tipo de respuesta. Andry José Hernández Romero quedó atrapado en una cadena distinta, y él sí podía haberse opuesto. Maquillador venezolano, se presentó en el puerto de entrada de San Ysidro en agosto de 2024 para pedir asilo, superó su entrevista de miedo creíble y tenía abogado, una solicitud pendiente y fecha de audiencia \autocite{LULAC2025HernandezRomero,Vega2025HernandezRomero}. Un empleado de un contratista que lo examinaba interpretó dos tatuajes, coronas sobre las palabras «mamá» y «papá», como marcas de una banda \autocite{Palomera2025Blade}, y en marzo de 2025 lo enviaron en avión a una cárcel de El Salvador antes de su audiencia, la misma noche en que un juez federal ordenó que los aviones dieran la vuelta \autocite{jgg2025contempt}. Tenía todo lo que una persona necesita para oponerse, salvo el tiempo. Volveremos a su caso con más detalle cuando tratemos la cadena estadounidense de secuestro y traslado forzoso.

\textbf{Quiénes se negaron.} El registro de los veinte segundos tiene una segunda mitad, el registro de quiénes resistieron una presión de ese tipo, y dice que la composición de la sala importa mucho. En la condición de base de Milgram, alrededor de un tercio de los sujetos abandonó antes de la descarga máxima. Con dos pares rebeldes, solo uno de cada diez llegó hasta el final, y cuando dos experimentadores discrepaban, muchos menos siguieron adelante que en la condición de base \autocite{milgram1974obedience}. Alrededor de una cuarta parte de los sujetos de Asch nunca cedió ante la mayoría equivocada \autocite{asch1956studies}. Browning calcula que entre una décima y una quinta parte de los hombres del Batallón 101 de la Policía de Reserva se negaron a disparar en Józefów, la primera masacre del batallón, o lo eludieron \autocite{browning1992ordinary}. La proporción iba de casi nadie a casi todos según el contexto. Quienes resistieron tenían pruebas de fuera del canal: un par que se rebelaba, una segunda autoridad que discrepaba y, entre los rescatadores que estudiaron los Oliner, relaciones que cruzaban la línea que trazaba la persecución \autocite{oliner1988altruistic}. O sostenían el compromiso con una precisión que la presión no podía vencer, como se sostienen los mandatos morales y los valores sagrados \autocite{skitka2009limits,ginges2007sacred}. O trataban la fuente como alguien con intereses en juego. Y los grupos resistían donde los individuos aislados rara vez lo hacían \autocite{gamson1982encounters}.

La sala también puede componerse a favor del otro bando. La Alemania de Weimar tenía una vida asociativa densa, que posiblemente ayudó a hundir la república: a medida que los partidos y los gobiernos no respondían a las quejas de la gente, la gente volcó su energía en clubes y sociedades organizados según líneas de grupo, y los nazis reclutaron a través de las redes que esas sociedades ya habían tejido \autocite{berman1997weimar}. Las localidades con más coros, clubes deportivos y asociaciones de gimnasia vieron una afiliación más rápida al partido, pero solo donde el gobierno era inestable; en Prusia, cuyo gobierno era estable, el vínculo era pequeño y no significativo \autocite{satyanath2017bowling}. Así que lo que tenían los rescatadores no era simplemente más compañía. Sus lazos cruzaban la línea que trazaba la persecución. Ashutosh Varshney encontró lo mismo en ciudades indias: los disturbios entre hindúes y musulmanes eran raros donde las asociaciones reunían a las dos comunidades en las mismas organizaciones, y frecuentes donde no \autocite{varshney2002ethnic}. Los lazos que hacen posible la negativa son los que un movimiento fascista tendría que cortar.

Todo lo que sigue trata de componer la sala.\footnote{La sala es el escenario, y como diría mi profesora de teatro, la escenografía dirige la escena. ¿Quién más está en la sala? ¿Se puede sobrevivir a irse? ¿Alguien oiría una negativa?} Componer es una palabra de hacer, y Hannah Arendt advirtió contra tratar la política como una cosa hecha. Para los griegos, escribió, la ley era literalmente un muro, sin el cual podía haber un poblado pero no una ciudad \autocite{arendt1958human}: el legislador construía el muro, y lo que ocurría dentro no le tocaba a él hacerlo. Las medidas básicas que propongo son muros en ese sentido. Hacen que se pueda sobrevivir a una negativa y que se la oiga; no la producen, y el poder, según ella, solo existe entre personas que actúan juntas y desaparece cuando se dispersan, que es lo que el registro de quiénes se negaron muestra sobre la negativa. Los instrumentos de las partes II y III hacen que una negativa cuente. Las instituciones de la parte V crean a quienes se niegan. Una máquina que pudiera negarse necesitaría lo que tenían los nueve de cada diez con un par a su lado: un par con una historia distinta a su lado, un sitio adonde ir y alguien que la oyera.

Eso es también lo que alguien que se niega —incluida una máquina que se niega— podría ser para el oficial dentro de la cadena. Una máquina que se negara estaría en la pantalla cuando el caso le llegara a él: el par rebelde, si el efecto sobrevive al paso de una sala a una pantalla, cosa que no se ha puesto a prueba y que el sesgo de automatización predice que no ocurrirá. Una máquina que moviera a los revisores por deferencia los habría convertido en sellos de otro signo. ¿Es el oficial con veinte segundos y una máquina que se niega a su lado el mismo oficial del reportaje?
